\VerdierPhysicalPageAnchor{0029}%
\VerdierSemanticAnchor{intro-7-duality-D4}
\noindent\textit{D4)} Lorsque $X$ est un fermé de $Y$ et que $f$ est
l'injection canonique, le foncteur $f^!$ est le foncteur dérivé total du
foncteur \og sections à support dans $X$\fg{}.

\VerdierSemanticAnchor{intro-7-traditional-duality}
Ce théorème, et ses variantes dans les autres contextes signalés plus haut,
contient les formulations traditionnelles de la dualité : théorème de
dualité de Poincaré dans le cas topologique, théorème de dualité de Serre pour
les faisceaux algébriques cohérents sur une variété projective et lisse. Mais
l'emploi du langage des catégories dérivées ne permet pas seulement de
démontrer avec plus de généralité des théorèmes auxquels on pouvait ne
s'intéresser que dans des cas particuliers simples. On ne sait par exemple,
à l'heure actuelle, démontrer le théorème de dualité sous la forme
traditionnelle (type dualité de Serre), pour les faisceaux algébriques
cohérents sur une variété propre et lisse non projective sur un corps
algébriquement clos, que par l'intermédiaire du théorème de dualité général.
Il en est de même pour la dualité en cohomologie étale, pour les variétés
projectives non singulières.

\VerdierSemanticAnchor{intro-8}
\noindent\textbf{8.}\quad
Les théorèmes de dualité dans les cas 1), 2) et 4) conduisent à des formules
de points fixes de Lefschetz, de type non classique, faisant intervenir des
faisceaux pouvant présenter des singularités [10], [11], [12]. Soient $X$ un
espace topologique compact de dimension cohomologique finie, $A$ un anneau
noethérien, $F$ un faisceau de $A$-modules sur $X$ possédant les propriétés
suivantes :

\VerdierSemanticAnchor{intro-8-lefschetz-hypotheses}
\begin{enumerate}
\renewcommand{\labelenumi}{\alph{enumi})}
\VerdierSemanticAnchor{intro-8-hypothesis-a}
\item En tout point de $X$, les fibres du faisceau $F$ sont des $A$-modules
de type fini et de dimension projective finie.

\VerdierSemanticAnchor{intro-8-hypothesis-b}
\item Une propriété de régularité au voisinage de tout point de $X$, analogue
aux propriétés décrites par Wilder [9].
\end{enumerate}

\VerdierSemanticAnchor{intro-8-endomorphism}
Soit $\Phi=(\Phi_X,\Phi_F)\colon(X,F)\longrightarrow(X,F)$ un endomorphisme
de l'espace topologique $X$ muni du faisceau $F$. L'endomorphisme $\Phi$
induit un endomorphisme de complexe de $A$-modules :
\VerdierSemanticAnchor{intro-8-rh-phi}
\[
\mathrm{RH}(\Phi)\colon
\mathrm{RH}(X,F)\longrightarrow\mathrm{RH}(X,F).
\]

\VerdierSemanticAnchor{intro-8-perfect-complex}
Sous les hypothèses a) et b), le complexe $\mathrm{RH}(X,F)$ est un
\emph{complexe parfait} de $A$-modules, \emph{i.e.} il est isomorphe, dans la
catégorie dérivée, à un complexe fini dont les composants sont des
$A$-modules projectifs de type fini.

\VerdierSemanticAnchor{intro-8-lefschetz-number}
On peut donc définir le nombre de Lefschetz $\chi(\mathrm{RH}(\Phi))$ comme
étant la somme alternée des traces des composants de $\mathrm{RH}(\Phi)$, qui
sont des endomorphismes de modules projectifs de type fini. Supposons
maintenant, pour fixer les idées, que les points fixes de $\Phi$ soient
isolés. On a alors l'égalité :
\VerdierSemanticAnchor{intro-8-lefschetz-formula}
\begingroup
\makeatletter
\tagsleft@true
\let\veqno\@@leqno
\makeatother
\begin{equation}
\chi\bigl(\mathrm{RH}(\Phi)\bigr)
=\sum_{P\ \mathrm{fixe}}\chi_P(\Phi)
\tag{Lef}
\end{equation}
\endgroup
